% Part: normal-modal-logic
% Chapter: axioms-systems
% Section: duals

\documentclass[../../../include/open-logic-section]{subfiles}

\begin{document}

\olfileid[hi]{nml}{prf}{dua}

\olsection{द्वैत \usetoken{P}{formula}}

\begin{defn}\ollabel{def:duals}
  !!{formula}ों \Ax{T}, \Ax{B}, \Ax{4} और \Ax{5} में से प्रत्येक का एक
  \emph{द्वैत} है, जिसे अधोलिखित हीरक से इस प्रकार निरूपित करते हैं:
  \begin{align}
    \tag{\Ax{T_\Diamond}} p & \lif \Diamond p\\
    \tag{\Ax{B_\Diamond}} \Diamond\Box p & \lif p\\
    \tag{\Ax{4_\Diamond}} \Diamond\Diamond p & \lif \Diamond p\\
    \tag{\Ax{5_\Diamond}} \Diamond\Box p & \lif \Box p
    \end{align}
\end{defn}

ऊपर का प्रत्येक द्वैत !!{formula}, संगत !!{formula} में $p$ के स्थान पर
$\lnot p$ रखकर, प्रतिधनात्मक करके, $\lnot\Box\lnot$ को $\Diamond$ से और
$\lnot\Diamond\lnot$ को~$\Box$ से बदलकर मिलता है। इस अर्थ में \Ax{D},
अर्थात् $\Box!A \lif \Diamond!A$, स्वयं अपना द्वैत है।

\begin{prop}\ollabel{prop:dualsys}
  \olref[nml][prf][dua]{def:duals} के प्रत्येक !!{formula}~$!A$ के लिए:
  $\Log{K}!A = \Log{K}!A_{\Diamond}$.
\end{prop}
\begin{proof}
  अभ्यास।
\end{proof}
\begin{prob}
  \olref[nml][prf][dua]{prop:dualsys} सिद्ध कीजिए।
\end{prob}

\end{document}
